\chapter{CONSIDERAÇÕES FINAIS}{}
Este trabalho apresentou o desenvolvimento de um sistema embarcado para
monitoramento de grandezas elétricas, voltado à aquisição de corrente e
tensão em tempo real e à posterior aplicação de técnicas de aprendizado
de máquina para previsão do consumo energético. Até o presente estágio,
foram concluídas as etapas de levantamento dos materiais, condicionamento
analógico dos sinais dos sensores SCT013-000 e ZMPT101B, desenvolvimento
do firmware embarcado no Raspberry Pi Pico 2 W e validação inicial do
\textit{hardware} desenvolvido.

Em relação aos objetivos específicos definidos, é
possível afirmar que os dois primeiros projetar um sistema embarcado
para aquisição de tensão e corrente, e realizar o monitoramento e
registro contínuo dessas grandezas foram atendidos, uma vez que o
sistema demonstrou capacidade de realizar leituras em regime permanente
e de registrar, de forma contínua, a transição para uma condição real
de distúrbio elétrico.

O terceiro objetivo, referente ao processamento dos sinais para
obtenção de parâmetros elétricos, foi parcialmente atendido: o sistema
calcula corretamente os valores RMS de corrente e tensão, restando
ainda a implementação do cálculo de potência ativa, reativa, aparente
e fator de potência, que depende da defasagem entre os sinais de
tensão e corrente de cada fase. O quarto objetivo, de armazenamento e
organização dos dados coletados, encontra-se atendido por meio do
registro em cartão microSD no formato CSV. Já o quinto objetivo,
relativo ao desenvolvimento de um modelo preditivo de consumo
energético baseado em aprendizado de máquina, ainda não foi iniciado
nesta etapa do trabalho, permanecendo como a principal atividade a ser
desenvolvida na continuidade do projeto.

Quanto à precisão das medições, o ensaio de validação indicou um erro relativo de
aproximadamente 4,3\% na leitura de corrente do sistema em relação ao
alicate amperímetro utilizado como referência. Embora esse valor esteja
dentro de uma faixa aceitável para um estágio inicial de calibração,
ele ainda é superior ao erro relatado por trabalhos correlatos
apresentados na revisão bibliográfica, como \citeonline{Ferraz2018}, o que indica a necessidade de um refinamento adicional do
processo de calibração de cada canal de corrente antes da validação
final do sistema.


\section{Próximas Etapas}


Diante do exposto, os resultados parciais obtidos até o momento indicam
que a arquitetura de \textit{hardware} e firmware desenvolvida é capaz
de atender aos requisitos básicos de aquisição, processamento e
registro de grandezas elétricas propostos neste trabalho. As próximas
etapas do desenvolvimento estão descritas no Quadro \ref{quadro:cronograma}.

\subsection{Cronograma de Atividades}

O Quadro~\ref{quadro:cronograma} apresenta o cronograma previsto para
a execução das próximas etapas do trabalho, entre julho e dezembro de
2026.

\begin{quadro}[htb]
\centering
\caption{Cronograma das próximas etapas do trabalho}
\label{quadro:cronograma}
\begin{tabular}{|p{7cm}|c|c|c|c|c|c|}
\hline
\textbf{Atividade} & \textbf{Jul} & \textbf{Ago} & \textbf{Set} & \textbf{Out} & \textbf{Nov} & \textbf{Dez} \\
\hline
Refinamento da calibração dos sensores de corrente & X & X & & & & \\
\hline
Cálculo de potência ativa, reativa e fator de potência & & X & X & & & \\
\hline
Critérios de validação temporal para eventos de tensão & & & X & X & & \\
\hline
Desenvolvimento da interface web & X & X & X & & & \\
\hline
Desenvolvimento da caixa (invólucro) do sistema & & X & X & & & \\
\hline
Teste do sistema com bateria de backup & & & X & X & & \\
\hline
Desenvolvimento e avaliação do modelo de aprendizado de máquina & & & & X & X & X \\
\hline
Redação final e revisão do TCC & & & & & X & X \\
\hline
\end{tabular}

\vspace{0.2cm}
\footnotesize{Fonte: Autoria Própria}
\end{quadro}